\chapter{결과: 수익, 안정성, 그리고 등급 안의 정보}

모든 결과는 2013--16년 발행 36개월 종결 대출(약 52만 건, 48개 발행월), 시점 학습, 투자자 수수료 1\% 차감 기준이다. 판정 기준은 각 실험 전에 정했다.

\section{출발점: LC 대출 자체는 국채보다 나았다}
2013--16년 36개월 대출을 \textbf{전부} 사면 발행월 묶음 IRR 5.73\%로 같은 기간 3년 국채(0.87\%)보다 높았고, 48개 발행월 \textbf{모두}에서 그랬다. 지난 우승팀은 단순 연환산과 최근 연도의 끝난 대출만 보는 설정에서 반대 결론(전부 사면 국채보다 못하다)을 냈다(2장). 그래서 우리의 질문은 ``LC에 투자할까''가 아니라 ``\textbf{무엇을 고르면 더 나은가}''가 된다.

\section{수익 축: 기대수익 순위가 가장 많이 번다}
\begin{table}[H]\centering\small
\begin{tabular}{@{}llr@{}}\toprule
비교 & 결과 & 근거\\\midrule
기대수익 순위 vs 무작위 & \textbf{+1.03\%p/년} (95\% 구간 0.95--1.11) & 25단계\\
기대수익 순위 vs 부도확률 순위 & +0.91\%p (0.73--1.10, 다 끝난 2013--15년) & 26단계\\
대출마다 같은 금액을 넣는 투자자 기준 & +1.12\%p (1.04--1.20) & 29단계\\
정답을 섞은 가짜 모델 & 0 근처 & 14단계\\\bottomrule
\end{tabular}
\caption{수익 축의 핵심 결과.}
\end{table}

\begin{table}[H]\centering\small
\begin{tabular}{@{}lrrl@{}}\toprule
모델 & AUC & 선별 가치 & 판정\\\midrule
LightGBM(기준) & 0.659 & 1.03\%p & 유지\\
XGBoost & 0.655 & 1.04\%p & 같음($p=0.41$)\\
랜덤포레스트 & \textbf{0.669} & 0.85\%p & AUC는 1위, 돈은 덜 범\\
Lasso + 피처 & 0.649 & 0.76\%p & 선형 중 최선\\
CatBoost & 0.651 & 0.68\%p & 2016년 $p$ 붕괴(평균 3.6\%)\\\bottomrule
\end{tabular}
\caption{모델 17종 중 주요 결과(23--25단계).}
\end{table}
두 가지를 읽을 수 있다. 첫째, \textbf{잘 맞히는 모델이 돈을 더 버는 것은 아니다}(AUC 1위 랜덤포레스트). 둘째, \textbf{1\%p는 모델의 산물이 아니다}. XGBoost가 같은 값을 냈고, 튜닝·\allowbreak{}새 피처·\allowbreak{}결측 처리로 바뀌지 않았다.

\paragraph{돈 단위로(NPV, 국채로 할인).} 승인 30\%에서 총 NPV는 기대수익 순 1억 2,800만 달러, 부도확률 순 1억 1,900만, 무작위 1억 100만 달러다. 9개 승인비율 모두 기대수익 순이 1위다. 다만 승인 50\% 이상에서는 부도확률 순과의 차이가 2--3\%뿐이다(33단계).

\section{등급 안의 정보: 이 프로젝트의 핵심 주장}
\subsection*{Brinson 분해}
기대수익 상위 50\%가 더 번 수익을 나누면 대부분이 \textbf{선별}(같은 등급 안에서 더 나은 대출)에서 나왔다. 세부등급 35개로 나눠도, 금리 구간·\allowbreak{}목적·\allowbreak{}지역·\allowbreak{}금액으로 묶어도, 60개월 대출에서도 같았다(14, 18단계).

\subsection*{등급 중립 롱숏}
\begin{table}[H]\centering\small
\begin{tabular}{@{}lrrr@{}}\toprule
신호 & 롱 $-$ 숏 & 90\% 구간 & 양수 달\\\midrule
기대수익 점수 & \textbf{+2.38\%p} & 2.08--2.73 & 48/48\\
최근 24개월 개설 계좌(적을수록) & +2.32\%p & 1.99--2.59 & 48/48\\
부도확률(낮을수록) & +2.01\%p & 1.72--2.34 & 48/48\\
연소득(높을수록) & +1.09\%p & 0.68--1.74 & 39/48\\
부채비율(낮을수록) & +0.95\%p & 0.76--1.24 & 44/48\\
FICO(높을수록) & +0.50\%p & 0.36--0.71 & 41/48\\
무작위(위약) & $-0.12$\%p & $-0.21$--$+0.03$ & 22/48\\\bottomrule
\end{tabular}
\caption{발행월·\allowbreak{}등급마다 신호 상위 30\% $-$ 하위 30\%의 IRR 차이(33단계). 실제로 가능한 롱온리(상위 30\% $-$ 그 등급 전체)는 기대수익 점수 +1.20\%p.}
\end{table}

\begin{keybox}[이 표가 말하는 것]
\begin{itemize}[leftmargin=*]
\item 등급과 시장 효과를 지워도, 같은 등급 안에서 고른 쪽이 \textbf{48개월 전부} 더 벌었다. 위약은 0이다.
\item \textbf{FICO는 약하다.} LC가 신용점수는 금리에 잘 반영했다.
\item \textbf{최근 계좌 개설 수 하나로 모델만큼의 신호}가 나온다. 짧은 기간에 대출·\allowbreak{}카드를 여러 개 연 사람은 돈이 급하거나 신용을 끌어다 쓰는 중일 가능성이 크고, LC는 이 행동을 금리에 덜 반영했다.
\end{itemize}
\Intuit{LC는 ``얼마나 신용이 좋은가''는 잘 쟀고, ``요즘 돈을 얼마나 찾아다니나''는 덜 쟀다.}
\end{keybox}

\subsection*{이 신호는 견고한가(34단계)}
\begin{table}[H]\centering\small
\begin{tabularx}{\linewidth}{@{}p{50mm}L@{}}\toprule
점검 & 결과\\\midrule
해마다 줄었나 & 기대수익 점수 1.94 \ensuremath{\to} 2.70 \ensuremath{\to} 1.81 \ensuremath{\to} 3.05\%p, 계좌 개설 1.79 \ensuremath{\to} 3.05\%p 근처로 줄지 않았다. 연소득만 2.08 \ensuremath{\to} 0.12\%p로 사라졌다\\
개인이 살 수 있던 대출(조각 대출)만 & +2.58\%p, 48/48개월. 조각 대출 비중은 72\% \ensuremath{\to} 27\%로 줄었다\\
60개월 대출(2013--14) & 계좌 개설 신호 +3.53\%p, 전 개월 양수\\
LC 가격(금리·\allowbreak{}등급) 없이 학습한 모델 & AUC 0.01 정도만 낮고, 등급 중립 롱숏 +2.17\%p\\\bottomrule
\end{tabularx}
\end{table}

\section{안정성 축: Sharpe로 보면 보수적 선별이 이긴다}
\begin{table}[H]\centering\small
\begin{tabular}{@{}lrrrr@{}}\toprule
승인 30\% & 포트폴리오 Sharpe & 평균 초과수익 & 건별 Sharpe(승인만) & 하위 5\% 평균\\\midrule
부도확률 낮은 순 & \textbf{3.24} & 1.66\%p & \textbf{0.26} & \textbf{$-20$\%}\\
금리 낮은 순 & 3.02 & 1.49\%p & 0.22 & $-21$\%\\
기대수익 순 & 2.46 & \textbf{1.92\%p} & 0.17 & $-32$\%\\
무작위 & 2.68 & 1.46\%p & 0.10 & $-34$\%\\\bottomrule
\end{tabular}
\caption{승인 30\%에서 두 축(30, 32, 33단계). 9개 승인비율 모두 같은 방향이고, 적게 승인할수록 차이가 크다.}
\end{table}
기대수익 순은 가장 많이 벌지만 꼬리 손실은 무작위와 거의 같다. 부도확률 순은 조금 덜 벌고 매달 훨씬 고르게 번다.

\subsection*{왜 그런가: 고위험 등급의 가격 실패}
\begin{table}[H]\centering\small
\begin{tabular}{@{}lrrr@{}}\toprule
등급(36개월) & 평균 초과수익 & Sharpe & 국채 이하 달\\\midrule
A & 4.0\%p & 5.2 & 0\\
B & 5.3\%p & 3.4 & 0\\
C & 5.3\%p & 2.3 & 0\\
D & 4.7\%p & 1.6 & 2\\
E & 3.3\%p & 0.6 & 17\\
F & 1.1\%p & 0.1 & 19\\
G & $-9.7$\%p & $-0.6$ & 6(8개월 중)\\\bottomrule
\end{tabular}
\caption{등급별 발행월 초과수익(33단계).}
\end{table}
Sharpe가 A에서 G로 내려가고, E--G는 평균 수익도 낮다. 36개월 고위험 대출은 \textbf{금리가 위험만큼 높지 않았다}. 기대수익 순은 금리가 높은 대출을 좋아하므로 이 영역의 출렁임을 떠안는다.

\begin{warnbox}[60개월은 다르다]
60개월(2013--14)에서는 평균 초과수익이 E(8.7\%p)까지 오르고 Sharpe는 B·\allowbreak{}C(6.7)가 가장 높았다. ``고위험 등급은 손해''는 36개월에만 맞는다. test 결과는 만기별로 따로 봐야 한다.
\end{warnbox}

\section{두 축의 화해}
등급 구성을 무작위와 같게 두고 등급 안에서만 부도확률 낮은 순으로 고르면, 무작위보다 수익은 크게 오르고(+0.34\%p) Sharpe는 조금만 오른다(2.68 \ensuremath{\to} 2.81). 여기에 안전한 등급 쪽으로 쏠리는 효과를 더하면(전체에서 부도확률 낮은 순) Sharpe는 크게 오르고(3.24) 수익은 내려간다.

\begin{keybox}[화해]
\textbf{수익은 같은 등급 안의 선별에서, 안정성은 안전한 등급으로의 배분에서 주로 나온다.} 우리 원래 주장(등급 안 선별)과 팀원 지적(보수적 선별의 Sharpe 우위)은 둘 다 맞았고, 서로 다른 축을 말하고 있었다.
\end{keybox}

\section{투명한 단순 규칙: 둘 다 잡을 수 있나}
``등급마다 최근 24개월 개설 계좌가 적은 대출을 승인''을 시험했다. 처음에는 수익·\allowbreak{}Sharpe 모두 1위였는데 \textbf{동점 버그}였다. 계좌 수가 정수라 경계값 동점이 모두 승인되어 실제 승인비율이 30\%가 아니라 39.5\%였다. 동점을 무작위로 끊어 다시 계산했다.

\begin{table}[H]\centering\small
\begin{tabular}{@{}lrr@{}}\toprule
승인 30\%(50\%) & 포트폴리오 수익 & 포트폴리오 Sharpe\\\midrule
단순 규칙 & 1.86\%p (2.88) & \textbf{3.72 (3.41)}\\
기대수익 순 & \textbf{1.92\%p (3.02)} & 2.46 (2.81)\\
부도확률 순 & 1.66\%p (2.83) & 3.24 (3.05)\\
무작위 & 1.45\%p (2.43) & 2.72 (2.69)\\\bottomrule
\end{tabular}
\caption{단순 규칙(34b단계).}
\end{table}
단순 규칙은 등급 구성을 바꾸지 않는데도 부도확률 순을 수익과 Sharpe \textbf{둘 다}에서 이긴다. ``안정성은 등급 배분에서만 나온다''가 아니었다. 무엇으로 고르느냐에 따라 등급 안 선별만으로 두 축을 다 얻을 수 있다. 다만 33단계에서 같은 데이터로 가장 강한 신호를 보고 고른 규칙이므로 \textbf{과대평가되었을 수 있다}. 그래서 test 규칙 v2에 가설 H7로 미리 고정했다(9장).

\section{투자자 목적별 메뉴}
\begin{table}[H]\centering\small
\begin{tabularx}{\linewidth}{@{}p{45mm}L@{}}\toprule
목적 & 규칙\\\midrule
가장 많이 벌기 & 기대수익 순\\
매달 고르게 벌기(과제의 Sharpe) & 부도확률 낮은 순, 또는 A·\allowbreak{}B등급\\
둘 다(후보, test에서 검증) & 등급 안 계좌 개설 적은 순\\
균형 & 등급 안 부도확률 낮은 순\\\bottomrule
\end{tabularx}
\end{table}
